\section{Mathematical tools}
\label{app:A}

In this appendix, we summarize the two tools used in Sections~\ref{sec:svm}
and~\ref{sec:disc}. Both of them are standard, but we give their intuitive meanings
and explain why they are necessary in this paper.

\subsection{The maximum theorem of Berge}
\label{app:A.1}

\begin{theorem}[Berge, 1963]
\label{thm:A.1}
Let $\Theta$ and $X$ be metric spaces, let $C:\Theta\rightrightarrows X$ be a
continuous correspondence with non-empty compact values, and let
$F:X\times\Theta\to\R$ be a continuous function. Put
\[
v(\vartheta)=\sup_{x\in C(\vartheta)}F(x,\vartheta),\qquad
M(\vartheta)=\arg\max_{x\in C(\vartheta)}F(x,\vartheta).
\]
Then $v$ is continuous and $M$ is upper hemicontinuous with non-empty compact
values. The upper hemicontinuity is written as
\begin{equation}
\vartheta_n\to\vartheta,\quad x_n\in M(\vartheta_n),\quad x_n\to x
\implies x\in M(\vartheta).
\label{eq:A.1}
\end{equation}
\end{theorem}

The content of the theorem is stated in one sentence as follows: the optimal value
moves continuously, while the optimal solution may jump, but the point to which it
jumps always belongs to the solution set of the limiting problem.

We give a minimal example. Let $X=\R$, $\Theta=\R$, $C(\vartheta)=[-1,1]$ and
$F(x,\vartheta)=\vartheta x$. Then $v(\vartheta)=|\vartheta|$ is continuous, while
\[
M(\vartheta)=\begin{cases}
\{1\} & \text{if }\vartheta>0,\\
[-1,1] & \text{if }\vartheta=0,\\
\{-1\} & \text{if }\vartheta<0,
\end{cases}
\]
so that the optimal solution jumps from $\{1\}$ to $\{-1\}$ at $\vartheta=0$.
Nevertheless~\eqref{eq:A.1} holds, because for $\vartheta_n\to 0$ we have
$x_n\in M(\vartheta_n)$ and indeed $x\in M(0)=[-1,1]$. That is, the upper
hemicontinuity states that the solution never appears at an irrelevant place, while
it allows the solution set to expand suddenly.

\begin{corollary}
\label{cor:A.1}
In the situation of Theorem~\ref{thm:A.1}, assume in addition that $M(\vartheta)$
is a singleton for each $\vartheta$. Then the map $\vartheta\mapsto M(\vartheta)$
is continuous.
\end{corollary}

\begin{proof}
Let $\vartheta_n\to\vartheta$ and $x_n=M(\vartheta_n)$. Take any subsequence of
$\{x_n\}$. From the compactness, it has a further subsequence converging to some
$x$. From~\eqref{eq:A.1} we have $x\in M(\vartheta)$, and from the uniqueness we
have $x=M(\vartheta)$. That is, every subsequence has a further subsequence
converging to the same limit $M(\vartheta)$, so that the whole sequence converges
to $M(\vartheta)$.
\end{proof}

We explain why this theorem is necessary. The solution of the SVM cannot be
written as an explicit function of $G$, because the set of the active constraints
changes with $G$. As long as the active set is fixed, the Karush--Kuhn--Tucker
conditions form a system of linear equations and one can obtain the
differentiability by the implicit function theorem. On the boundary at which the
active set switches, however, the differentiability is lost in general. What we
need in this paper is not the differentiability but only the continuity, and
Theorem~\ref{thm:A.1} gives it from the continuity of the objective function and
the feasible correspondence alone.

In order to apply the theorem, we have to supply the two conditions. As for the
continuity of the correspondence, we derive the a priori bound $\|\alpha\|\le K_0$
from the strict concavity and restrict the feasible set to a compact set which
does not depend on $G$. A constant correspondence is trivially continuous, so that
the most delicate condition of the theorem is met without any effort. As for the
uniqueness, we use Lemma~\ref{lem:4}.

\subsection{Fatou's lemma}
\label{app:A.2}

\begin{theorem}[Fatou's lemma]
\label{thm:A.2}
Let $\{f_n\}$ be a sequence of non-negative measurable functions on a measure
space $(X,\mathcal{F},\nu)$. Then
\begin{equation}
\int\liminf_{n\to\infty}f_n\,d\nu
\le\liminf_{n\to\infty}\int f_n\,d\nu.
\label{eq:A.2}
\end{equation}
\end{theorem}

In the language of probability, for non-negative random variables $\{X_n\}$ it
holds that $E\{\liminf X_n\}\le\liminf E\{X_n\}$.

The content of the lemma is stated as follows: a non-negative mass may escape in
the limiting procedure, but it never emerges.

We give the standard example in order not to mistake the direction of the
inequality. Let $X=[0,1]$ with the Lebesgue measure and $f_n=n\cdot I_{[0,1/n]}$.
For each $x>0$ we have $f_n(x)=0$ for sufficiently large $n$, so that
$\liminf f_n=0$ almost everywhere. On the other hand, $\int f_n=1$ for every $n$.
Hence~\eqref{eq:A.2} is a strict inequality. The mass of one escaped through the
pointwise convergence as a thin spike of height $n$ and width $1/n$. We note that
the non-negativity is necessary, because for $g_n=-f_n$ we have
$\int\liminf g_n=0$ and $\liminf\int g_n=-1$, so that~\eqref{eq:A.2} does not hold.

We explain why this lemma is necessary. In Corollary~\ref{cor:1}, what we have is
the convergence in distribution $e(1)\wto e^*(1)$ only, which does not imply the
convergence of the expectation in general. On the other hand, what we need is the
lower bound $E\{e^*(1)\}\le\liminf E\{e(1)\}$, which is exactly the direction of
the inequality in~\eqref{eq:A.2}. That is, in order to show the inconsistency, the
lower bound suffices and Fatou's lemma is enough.

Since~\eqref{eq:A.2} is a statement about the almost sure convergence, we use the
following bridge.

\begin{lemma}
\label{lem:7}
If $X_d\ge 0$ and $X_d\wto X$, then $E\{X\}\le\liminf_d E\{X_d\}$.
\end{lemma}

\begin{proof}
From the Skorokhod representation theorem \citep{billingsley1999}, there exist
random variables $Y_d$ and $Y$ on a common probability space such that $Y_d$ and
$X_d$ have the same distribution, $Y$ and $X$ have the same distribution, and
$Y_d\to Y$ almost surely. Applying Theorem~\ref{thm:A.2} to $\{Y_d\}$
\citep{durrett2019}, we obtain $E\{Y\}\le\liminf E\{Y_d\}$. Since the
expectation depends only on the distribution, we obtain the result.
\end{proof}

\begin{remark}
\label{rem:A.1}
In our case, we actually have the equality. Since $e(1)$ is a probability, it
holds that $0\le e(1)\le 1$, that is, the sequence is uniformly bounded. Applying
the Portmanteau theorem to the identity function on $[0,1]$, we obtain
$E\{e(1)\}\to E\{e^*(1)\}$. We use Fatou's lemma because the lower bound is what
we need and because it does not require the uniform boundedness. When one measures
the error rate in another scale which is not bounded, the lower bound given by
Fatou's lemma becomes essential.
\end{remark}
